\chapter*{TÓM TẮT}
Việc ôn luyện cho các kỳ thi tiếng Anh chuẩn hóa như IELTS, TOEIC, TOEFL và SAT đòi hỏi người học không chỉ tiếp cận được học liệu phù hợp mà còn cần một môi trường luyện tập có cấu trúc, phản hồi kịp thời và khả năng theo dõi tiến độ trong thời gian dài. Tuy nhiên, nhiều nền tảng hiện nay vẫn chủ yếu tập trung vào ngân hàng đề hoặc bài luyện tĩnh, trong khi các chức năng hỗ trợ cho kỹ năng tạo sinh như Speaking và Writing, cũng như khả năng kết nối giữa tự học với phản hồi chuyên sâu, vẫn còn hạn chế.

Báo cáo này trình bày quá trình thiết kế và hiện thực một nền tảng hỗ trợ ôn luyện tiếng Anh ứng dụng trí tuệ nhân tạo, tích hợp nhiều nhóm chức năng trong cùng một hệ thống thống nhất. Ở mức nền tảng, hệ thống cung cấp các mô-đun phục vụ học tập và quản trị như quản lý người dùng, tìm kiếm và thực hiện bài thi thử, theo dõi lịch sử làm bài, quản lý mục tiêu học tập, blog, báo cáo và thông báo. Về công nghệ, hệ thống được xây dựng theo kiến trúc web hiện đại với Next.js cho giao diện người dùng, NestJS cho các dịch vụ backend, kết hợp PostgreSQL, Minio và các dịch vụ AI chuyên biệt nhằm bảo đảm tính mở rộng và sự tách biệt giữa các miền chức năng.

Điểm nhấn của đề tài nằm ở hai mô-đun AI cho Speaking và Writing. Đối với Speaking, nhóm xây dựng một nền tảng luyện nói riêng, nơi người học có thể luyện phản xạ nói, nhận phản hồi và kết nối với giáo viên hoặc mentor trong cùng một chu trình học tập kết hợp giữa AI và con người. Hệ thống sử dụng mô hình ngôn ngữ lớn để sinh phản hồi phù hợp với ngữ cảnh học tập. Đối với Writing, nhóm phát triển mô-đun chấm và phản hồi bài viết dựa trên mô hình Phi-3 được fine-tune trên dữ liệu IELTS Writing, sau đó triển khai thông qua Writing service và Hugging Face Inference Endpoint để hỗ trợ suy luận ổn định.

Nền tảng này góp phần khắc phục những hạn chế của các giải pháp ôn thi tiếng Anh hiện có bằng cách cung cấp môi trường luyện tập thuận tiện hơn, mô phỏng sát thực tế hơn và cơ chế phân tích kết quả dựa trên AI. Quá trình triển khai cho thấy khả năng kết hợp hiệu quả giữa phát triển web hiện đại và trí tuệ nhân tạo để tạo nên một giải pháp công nghệ giáo dục có tính thực tiễn, khả năng mở rộng và định hướng lấy người học làm trung tâm.